\scanpage{5}
effet sinon il existerait un homomorphisme non nul d'un $M\handwritten{\in}\underline G$ dans $F/F'$,
et l'image inverse \inserted{dans $F$} de son image \deleted{dans} contredirait l'hypothèse maximale
pour $F'$.

\uline{Exemple} 2.4. Soient $A$ un anneau comm noeth, $I$ un idéal de $A$, $X$ le spectr\unsure{e}
de $A$ et $Y$ la partie fermée de $X$ définie par $I$. Soit $\underline C$ la catégorie des
$A$-modules \inserted{de type fini} dont tout élément est annulé par une puissance de $I$, ou la catégorie équivalente des Modules \deleted{quasi-}cohérents sur $X$ dont le support est
contenu dans $Y$. Les objets de $\underline C$ sont noethériens. Si $\mathbb M\handwritten{\in}\mathrm{Ob}\,\underline C$, alors on
voit tout de suite que $\mathbb M$ \uline{est un générateur de} $\underline C$ \uline{sss} $\operatorname{supp}\mathbb M=Y$. C'est
nécessaire, car si on avait $\operatorname{supp}\mathbb M\handwritten{\ne}Y$, il existerait une composante
irréductible $Y_i$ de $Y$ non contenue dans\deleted{\unsure{s}} $\operatorname{supp}M$, et prenant $F=\underline O_{Y_i}$, avec
$Y_i$ muni de la structure réduite induite, on trouve que $\operatorname{Hom}(M,F)=0$, $F\handwritten{\ne}0$.
C'est suffisant, car si \deleted{supp} $F$ est \deleted{\unsure{selement|seulement}} quasi-cohérent, de support
$\handwritten{\subset}Y$, on a $\operatorname{Ass}\operatorname{Hom}(M,F)=\operatorname{supp}M\handwritten{\cap}\operatorname{Ass}F=\operatorname{Ass}F$, donc $\operatorname{Hom}(M,F)=0$ implique
$\operatorname{Ass}F=\handwritten{\varnothing}$ d'où $F=0$. \deleted{\unsure{(}}(Donc $M$ est un générateur également dans la catégorie de tous les m\deleted{\unsure{d}}odules \uline{quasi}-cohérents $F$ de support $\handwritten{\subset}Y$). \deleted{Par suite,
si $T$ pour tester la nullité de} Le même argument montre que si on a une
famille $(M_i)_{i\handwritten{\in}I}$ d'objets de $\underline C$, alors c'est un système de générateurs de $\underline C$
sss la \deleted{\unsure{re}}réunion des $\operatorname{supp}M_i$ est $Y$ (commencer par ça). Pour une famille
satisfaisant cette condition, l'hypothèse de 2.1. est vé\deleted{\unsure{r}}rifée, et par
suite, pour vérifier la nullité de foncteurs comme dans 2.1. ou 2.2.,
il suffit de tester par les $M_i$, \deleted{ou de tester \unsure{i}} en particulier il suffit
de tester par \uline{un} $M$ de support $Y$.

\uline{Exemple} 2.5. On peut généraliser 2.4. en prenant un schéma \deleted{\inserted{quasi-compact}}\deleted{$X$ muni d'un}
noethérien $X$, muni d'un faisceau ample noté $\underline O_X(1)$, \deleted{Soit $(M_i)$} et d'une
partie fermée $Y$, $\underline C$ désignant encore la catégorie des faisceaux cohérents
sur $X$ à support $\handwritten{\subset}Y$. Les objets de $\underline C$ sont encore noethériens, et si
$(M_i)_{i\handwritten{\in}I}$ est une famille d'objets de $\underline C$, telle que \deleted{$\operatorname{supp}M_{\unsure{Y|i}}\unsure{\subset|=}$}
\[
\handwritten{\bigcup_i}\operatorname{supp}M_i=Y\quad,
\]
(il aurait fallu déployer la formule é\handwritten{g}alement dans 2.4.) , alors les
$M_i(\deleted{\unsure{+}}\unsure{-|\pm}n)$, avec $n\handwritten{>}N$ ($N$ entier donné), forment un système de générateurs
de $\underline C$ (et même de la catégorie des Modules qu coh de support $Y$). En effet
si $F$ est un Module qu coh de support $Y$, et $F\handwritten{\ne}0$, alors \deleted{\unsure{l'arg}} on a vu
dans 2.4. que le faisceau $\underline{\operatorname{Hom}}(M_{\deleted{\unsure{n}}i},F)$ est non nul pour $i$ convenable (se placer dans un ouvert affine où $F$ est $\handwritten{\ne}0$), donc il existe des sections non
nulles de $\underline{\operatorname{Hom}}(M_i,F)(n)$ pour $n$ grand, donc des\handwritten{\unsure{/}}homomorphismes non nuls de
$M_i(-n)$ dans $F$ pour $n$ grand. -- On peut donc utiliser les $M_i(-n)$, $n\handwritten{>}N$,
% archive sheet number: 4

